\section{Coverage ledger, conclusions, and further research}\label{sec:conclusion}

\subsection{Coverage of the source article}

The accompanying \code{THEOREM\_LEDGER.md} contains 94 numbered mathematical
items extracted from the consulted revision of the article.  The coverage is
organized as follows.

\begin{longtable}{@{}p{0.22\textwidth}p{0.68\textwidth}@{}}
\toprule
Article region & Constructive treatment \\
\midrule
\endhead
Finite definition and examples & Exact natural-weight/rational-law
formalization; die and roulette computations; equiprobable arithmetic mean;
vector and matrix expectation. \\
Countable definition & Absolute-tail regular sums; effective reindexing;
conditional rearrangement counterexample; logarithmic law; extended
expectations; St Petersburg divergence; undefined two-sided example. \\
Density and general integral & Effective simple-law names, compact quadrature,
improper tails, density agreement, LOTUS, Stieltjes integration, and layer
cake.  Existence-only density extraction is replaced by a witness. \\
Common distributions & Bernoulli and binomial exactly; Poisson, geometric,
uniform, exponential, normal, Pareto, and Cauchy with explicit series or
integration moduli. \\
Basic properties & Nonnegativity, finite linearity, monotonicity,
nondegeneracy, almost-sure congruence, constants, absolute-value bound,
indicators, CDF/tail formulas, nonmultiplicativity, independence, and
multidimensional LOTUS. \\
Inequalities & Markov, Chebyshev, certified Jensen, Lyapunov, Holder,
Cauchy--Schwarz, and Minkowski. \\
Convergence & Effective monotone convergence, nonnegative countable additivity,
Fatou, dominated convergence, uniform integrability, plus exact counterexamples
when hypotheses are omitted. \\
Characteristic functions & Derivative-at-zero estimate, finite cyclic Fourier
expectation identity, and continuous effective inversion under explicit
Fubini/tail hypotheses. \\
Applications & Sample means, residual squares, moment determinacy on bounded
support, expected utility, indicators and Monte Carlo, center of mass, beta
balance, bounded Wald identity, variance, and finite-dimensional quantum
expectation. \\
\bottomrule
\end{longtable}

Historical claims, naming conventions, and etymology were not treated as
mathematical proof obligations.  The report's aim is not to reconstruct the
article's prose but to identify its formal mathematical core.

\subsection{Main constructive conclusion}

The expected-value operator is unusually well suited to \(\ipra\).  Its most
important algebraic laws are consequences of one primitive-recursive bounded
sum.  Even many probabilistic inequalities are pointwise finite inequalities
followed by that same sum.  The logical strength lies not in measure theory but
in the representation of the data.

The exact finite theory therefore needs no completed real line, sigma-algebra,
or countable additivity.  It already includes enough probability for
finite experiments, exact randomized algorithms, finite statistical models,
bounded stopping times, finite-dimensional quantum systems, and explicit
finite-error laws of large numbers.

For infinite objects, the correct constructive currency is a modulus.  An
absolute-tail modulus makes a countable expectation order independent.  An
\(L^1\) modulus makes expectations of simple approximants Cauchy.  A quadrature
modulus and an improper-tail modulus turn an integral into a regular real.  A
uniform-integrability modulus licenses limit interchange.  Once these data are
present, the proofs again reduce to the finite kernel.

\subsection{Corrections and cautions}

Three points deserve emphasis.

\begin{enumerate}
\item A conditionally convergent weighted series is not an expectation merely
  because one enumeration converges.  The outcomes have no canonical order;
  absolute convergence or a positive/negative-part definition is essential.
\item A symmetric principal value is not an expectation.  The Cauchy law has
  zero symmetric truncations but different asymmetric limits, while both
  absolute parts diverge.
\item The die/Chebyshev value \(8/15\) concerns the strict event
  \(1<X<6\), not the inclusive event that a die lies in its support.  The
  exact strict-event probability is \(2/3\), and the inclusive probability is
  one.
\end{enumerate}

\subsection{What the packet does not yet prove in Lean}

The human proofs are designed for compilation, but the packet stops before the
final kernel-checked endpoint.  In particular:
\begin{itemize}
\item the ambient Lean theorem-package records are not instantiated;
\item the natural decoders and common-denominator compiler are not implemented;
\item no concrete semantic checker has yet been compiled to \code{PR};
\item no new \code{Prf PRAAx} derivation is present;
\item the analytic modules are interfaces rather than completed interval
  arithmetic.
\end{itemize}
These limitations are explicit so that standard-model mathematics is not
mistaken for an internal proof.

\subsection{Further research questions}

\begin{enumerate}
\item \textbf{Minimal finite expectation kernel.}
  What is the smallest extension of the current ProveIt bounded-sum library
  that makes linearity, finite LOTUS, and Markov one- or two-line internal
  derivations?  A dedicated weighted-sum abstraction may reduce proof size
  substantially without increasing trusted strength.

\item \textbf{Canonical denominator compiler.}
  Can a verified compiler transform ambient \code{Rat} finite-sum theorems
  into signed-natural, denominator-cleared \code{PRAAx} goals automatically,
  while exposing a compact proof certificate checked by the existing kernel?

\item \textbf{Probability-specific reflection.}
  The Lagrange identity, variance expansion, and binomial expectation are
  polynomial identities after finite sums are normalized.  A reflection
  tactic that combines \code{MultiSums} normalization with \code{deriveNat}
  could cover a broad class of probability proofs.

\item \textbf{First-crossing library.}
  Kolmogorov's inequality, optional-stopping arguments, and bounded martingale
  theorems share a primitive-recursive first-hit construction.  A reusable
  first-crossing module may unlock many finitary stochastic-process results in
  \(\ipra\).

\item \textbf{Optimal effective hypotheses.}
  The analytic theorems here use convenient \(L^1\), tail, and probability
  moduli.  Which of these can be weakened while retaining a primitive-recursive
  output modulus?  In particular, can effective almost-sure convergence be
  represented more economically than convergence-in-probability plus uniform
  integrability?

\item \textbf{Constructive Radon--Nikodym fragments.}
  For which finitely generated algebras, computable compact spaces, or bounded
  density classes can a density be extracted in \(\ipra\) from a quantitative
  absolute-continuity modulus?  The general classical converse is too weakly
  specified, but important restricted forms may be fully effective.

\item \textbf{Exact cyclotomic Fourier layer.}
  Is a small verified cyclotomic-ring library simpler than regular sine/cosine
  names for finite Fourier inversion?  Such a layer would also benefit finite
  characteristic functions and discrete harmonic analysis.

\item \textbf{Formal expected utility and games.}
  The finite maximization theorem is immediate.  A natural next target is a
  verified constructive account of the problem of points, finite zero-sum
  games with rational mixed strategies, and backward induction, all phrased as
  expectation theorems.

\item \textbf{Moment problems at bounded complexity.}
  Finite support is exactly determined by finitely many moments.  What useful
  effective determinacy results hold for distributions with a supplied tail
  envelope, finite mixture models, or compact algebraic supports, within
  feasible \(\ipra\) proof size?

\item \textbf{Compression of analytic certificates.}
  The regular-real approach is transparent but potentially verbose.  Can
  Taylor, quadrature, and tail certificates be compressed into a single
  proof-producing evaluator whose correctness is established once in
  \(\mathrm{PRAAx}\)?
\end{enumerate}

\subsection{Final assessment}

A substantial portion of elementary expectation theory is not merely
constructively acceptable; it is primitive-recursively explicit.  The finite
kernel contains the conceptual heart of the subject, while the infinite theory
is a disciplined closure under named limits.  The remaining challenge is
proof engineering: implementing the encoders and turning the arguments already
reduced to finite algebra into compact checked \code{PRAAx} derivations.
